\section{Multiplication determinants without large matrices}
\label{sec:algebra}

The geometric constructions below produce polynomial systems whose
leading terms are pure powers.  Their quotient algebras have explicit
monomial bases, but those bases can be exponentially large.  The goal
of this section is to compute characteristic polynomials of
multiplication in these algebras as controlled families, without
constructing the multiplication matrices.  We first express a trace
as a coefficient, and then recover characteristic coefficients from
traces by fixed-depth finite-field computations.

\subsection{A basis and a coefficient functional}

Let $S$ be a commutative ring with identity, let $n\ge1$ and $\ell\ge2$,
and suppose that $r_1,\ldots,r_n\in S[x_1,\ldots,x_n]$ satisfy
$\deg_x r_i<\ell$.  Set
\begin{equation}\label{eq:quotient-algebra}
 f_i=x_i^\ell-r_i,\qquad
 \mathcal A=S[x_1,\ldots,x_n]/(f_1,\ldots,f_n),\qquad
 D=\ell^n,\qquad \tau=(\ell-1,\ldots,\ell-1).
\end{equation}
For a multi-index $a$, write $|a|=\sum_i a_i$ and
$x^a=\prod_i x_i^{a_i}$.  We will show that the $D$ monomials
$x^a$ with $0\le a_i<\ell$ form an $S$-basis of $\mathcal A$.

We use formal Laurent sums whose total degrees are bounded above and
whose support in each total degree is finite.  Addition and
multiplication are well defined: for any prescribed total degree,
only finitely many pairs of degrees and finitely many terms can
contribute.  Since $r_i x_i^{-\ell}$ has total degree at most $-1$,
its geometric series is such a Laurent sum.  Define the $S$-linear
functional
\begin{equation}\label{eq:laurent-functional}
 \mathcal L(v)=
 [x_1^{-1}\cdots x_n^{-1}]
 v(x)\prod_{i=1}^n
 \left(x_i^{-\ell}\sum_{j\ge0}
       (r_i(x)x_i^{-\ell})^j\right).
\end{equation}
This is a formal coefficient extraction, requiring neither convergence
nor division in $S$.  It is a concrete residue formula of the type
studied in \cite{CattaniDickensteinSturmfels1996}; we prove the facts
needed here directly over $S$.

\begin{lemma}\label{lem:monomial-basis}
The functional $\mathcal L$ vanishes on $(f_1,\ldots,f_n)$ and hence
descends to $\mathcal A$.  If $\deg_x v\le n(\ell-1)$, then
\begin{equation}\label{eq:low-degree-functional}
 \mathcal L(v)=[x^\tau]v.
\end{equation}
The algebra $\mathcal A$ is a free $S$-module with basis
$\{x^a:0\le a_i<\ell\}$.
\end{lemma}

\begin{proof}
Multiplying the expression in Equation~\eqref{eq:laurent-functional}
by $f_i=x_i^\ell-r_i$ cancels its $i$th geometric factor.  In the
remaining expression, every exponent of $x_i$ is nonnegative:
the other denominators are powers of the other variables, while all
numerators are polynomials.  Its coefficient with $x_i$-exponent $-1$
is zero.  Thus $\mathcal L(vf_i)=0$ for every polynomial $v$.

Expanding all the geometric factors gives the finite sum
\begin{equation}\label{eq:laurent-coefficient-sum}
 \mathcal L(v)=\sum_{j\in\mathbb Z_{\ge0}^n}
 [x^{\tau+\ell j}]\left(v\prod_{i=1}^n r_i^{j_i}\right).
\end{equation}
The target degree in a summand is $n(\ell-1)+\ell|j|$, whereas the
numerator has degree at most $\deg_x v+(\ell-1)|j|$.  A nonzero
summand therefore requires
\begin{equation}\label{eq:residue-truncation}
 |j|\le\deg_x v-n(\ell-1).
\end{equation}
This proves finiteness and Equation~\eqref{eq:low-degree-functional}.

The relations $x_i^\ell=r_i$ express every monomial as a linear
combination of the proposed basis elements: each replacement strictly
decreases total $x$-degree.  For independence, consider
\[
 P_{a,b}=\mathcal L(x^a x^{\tau-b}),
 \qquad a,b\in\{0,\ldots,\ell-1\}^n.
\]
Equation~\eqref{eq:low-degree-functional} gives $P_{a,b}=0$ if
$|a|<|b|$, and $P_{a,b}=\delta_{a,b}$ if $|a|=|b|$.
Ordering both index sets by total degree makes $P$ block triangular
with identity diagonal blocks.  It is therefore invertible over $S$.
Pairing any linear relation among the proposed basis elements against
all $x^{\tau-b}$ forces all its coefficients to be zero.
\end{proof}

\subsection{A trace formula and its spectral consequences}

Let $M_v$ denote the matrix of multiplication by the class of
$v\in S[x]$ in this basis, and put
\begin{equation}\label{eq:jacobian}
 J(x)=\det\bigl(\partial_j f_i(x)\bigr)_{1\le i,j\le n}.
\end{equation}
The following classical trace--residue identity is associated with
Scheja and Storch \cite{SchejaStorch1975}; see also
\cite[Author version, Equation~(4.7)]{CattaniDickensteinSturmfels1996}.
Its proof uses the dual-basis
property of a Bezoutian, a construction developed further in
\cite[Section~3.3]{ElkadiMourrain2005}.  The explicit proof matters here because
we need arbitrary coefficient rings and specialization at multiple
zeros.

\begin{proposition}[Trace formula]\label{prop:trace}
For the algebra in Equation~\eqref{eq:quotient-algebra}, over every
commutative ring $S$ with identity,
\begin{equation}\label{eq:trace-formula}
 \operatorname{tr}(M_v)=\mathcal L(vJ)
 \qquad(v\in S[x]).
\end{equation}
\end{proposition}

\begin{proof}
Introduce another set of variables $x'=(x'_1,\ldots,x'_n)$.
Changing one variable at a time gives polynomial divided differences
$W_{ij}(x,x')$ such that
\[
 f_i(x)-f_i(x')=\sum_{j=1}^n
 W_{ij}(x,x')(x_j-x'_j).
\]
Each $W_{ij}$ has combined $(x,x')$-degree at most $\ell-1$.
Let $W=\det(W_{ij})$.  In $\mathcal A\otimes_S\mathcal A$, the
left sides vanish, so multiplication by the adjugate matrix gives
\begin{equation}\label{eq:diagonal-annihilation}
 W(x,x')(x_j-x'_j)=0\qquad(1\le j\le n).
\end{equation}

Apply $\mathcal L$ in the second factor, writing this map as
$\mathcal L_{x'}$.  Since $\deg_{x'}W\le n(\ell-1)$,
Lemma~\ref{lem:monomial-basis} gives
$\mathcal L_{x'}(W)=[(x')^\tau]W$.  This coefficient is $1$.
Indeed, the part of combined degree $\ell-1$ in the divided-difference
matrix is diagonal, with $i$th diagonal entry
\[
 x_i^{\ell-1}+x_i^{\ell-2}x'_i+\cdots+(x'_i)^{\ell-1}.
\]
The contributions from the $r_i$ have combined degree at most
$\ell-2$.  To obtain the $(x')^\tau$ coefficient, whose total
degree is $n(\ell-1)$, one must use the last term of every displayed
diagonal entry.  The target exhausts the combined-degree bound, so
no positive $x$-degree can occur in this coefficient.

Equation~\eqref{eq:diagonal-annihilation} now yields the reproducing
identity
\begin{equation}\label{eq:reproducing-identity}
 \mathcal L_{x'}\bigl(W(x,x')v(x')\bigr)=v(x)
 \quad\text{in }\mathcal A:
\end{equation}
within a product with $W$, each $x'_j$ can be replaced by $x_j$.
Expand $W=\sum_a x^a w_a(x')$ in the first monomial basis of the
tensor product.  Equation~\eqref{eq:reproducing-identity} says that
the $x^a$-coordinate of $v$ is $\mathcal L(w_a v)$.  In particular,
the $a$th diagonal entry of $M_v$ is $\mathcal L(w_a v x^a)$.
The multiplication map $\mathcal A\otimes_S\mathcal A\to\mathcal A$
identifies $\sum_a x^a w_a(x)$ with $W(x,x)$, so
\[
 \operatorname{tr}(M_v)
 =\mathcal L\!\left(v\sum_a x^a w_a(x)\right)
 =\mathcal L\bigl(vW(x,x)\bigr).
\]
On the diagonal $x'=x$, divided differences become derivatives;
hence $W(x,x)=J(x)$, proving the formula.
\end{proof}

Both the basis and the trace formula commute with every change of
coefficient ring.  In particular, parameters may be specialized and
coefficients reduced modulo a prime without assuming that the common
zeros are distinct.  The following elementary fact explains how the
characteristic polynomials will encode real critical values.

\begin{lemma}\label{lem:real-eigenline}
Suppose that the coefficients of the relations in
Equation~\eqref{eq:quotient-algebra} and of $v$ lie in $\mathbb C$,
and write $Q(y)=\det(yI-M_v)$.  Every common zero
$\xi\in\mathbb C^n$ of the $f_i$ satisfies $Q(v(\xi))=0$.
If all coefficients are real and $\eta$ is a simple real root of $Q$,
there is a common zero $\xi\in\mathbb R^n$ with $v(\xi)=\eta$.
Moreover, for any polynomial $g$ over $\mathbb C$,
\begin{equation}\label{eq:composition-determinant}
 M_{g(v)}=g(M_v),\qquad
 \det M_{g(v)}=\prod_{Q(\alpha)=0}g(\alpha),
\end{equation}
where the roots in the product are counted with multiplicity.
\end{lemma}

\begin{proof}
Evaluation at a common zero $\xi$ is a nonzero linear functional on
$\mathcal A$, since it sends $1$ to $1$.  It satisfies
$\operatorname{ev}_\xi(vw)=v(\xi)\operatorname{ev}_\xi(w)$,
so it is a left eigenvector of $M_v$ with eigenvalue $v(\xi)$.

For real coefficients, a simple real eigenvalue $\eta$ of the real
matrix $M_v$ has a one-dimensional real eigenspace.  Each coordinate
multiplication matrix $M_{x_i}$ commutes with $M_v$ and therefore
preserves this line.  Let $\xi_i\in\mathbb R$ be its scalar action
there.  Restricting the matrix identities
\[
 f_i(M_{x_1},\ldots,M_{x_n})=0,\qquad
 v(M_{x_1},\ldots,M_{x_n})=M_v
\]
to the line gives $f_i(\xi)=0$ and $v(\xi)=\eta$.
Finally, multiplication respects sums and products, giving the first
identity in Equation~\eqref{eq:composition-determinant}; triangularizing
$M_v$ over $\mathbb C$ gives the second.
\end{proof}

\subsection{Characteristic coefficients from finite-field evaluations}

We now turn the trace formula into a uniform computation.  All the
large sums below are indexed by tuples of polynomial bit length.
Each entire tuple is a single counting index; the number of its
coordinates does not increase the number of counting levels.
We state the evaluation bound separately, so that later applications
with elementary evaluators can retain their sharper level count.

\begin{proposition}[Indexed characteristic coefficients]
\label{prop:char-coeff}
Consider fixed families $r_1,\ldots,r_n,v\in\mathbb Z[x]$, indexed
by a word of length $N$.  Suppose that $n\ge1$ is polynomially bounded
and polynomial-time computable, that $\ell\ge2$ and $\Omega\ge0$
are short polynomial-time computable integers, and that
$\deg_x r_i<\ell$ and $\deg_x v\le\Omega$.  Put
\[
 f_i=x_i^\ell-r_i,\qquad
 J=\det(\partial_jf_i),\qquad
 D=\ell^n,\qquad L_0=D\Omega+n(\ell-1).
\]
For $a\ge0$, suppose that $v$, the indexed $r_i$, and $J$ can be
evaluated uniformly in $\Flevel{a}$ over prime fields with
\begin{equation}\label{eq:char-prime-bound}
 p>n(\ell-1)+\ell L_0+D^2+10.
\end{equation}
Additional short, polynomial-time computable lower cutoffs on $p$
are allowed.  Then, for every $0\le d\le D$,
\[
 [y^d]\det(yI-M_v)\pmod p
\]
is uniformly computable in $\Flevel{a+8}$ in this prime range.
All running-time and query-length bounds are polynomial in the
length of the input, including $p$ and the coefficient index $d$.
\end{proposition}

\begin{proof}
First compute the trace residues $t_k=\operatorname{tr}(M_v^k)\bmod p$
for short indices $1\le k\le D$.  By Proposition~\ref{prop:trace},
$M_v^k=M_{v^k}$, and Equation~\eqref{eq:laurent-coefficient-sum},
\begin{equation}\label{eq:trace-extraction}
 t_k=\sum_{\substack{j\in\mathbb Z_{\ge0}^n\\|j|\le L_0}}
 [x^{\tau+\ell j}]
 \left(v^kJ\prod_{i=1}^n r_i^{j_i}\right)\pmod p.
\end{equation}
Indeed, $\deg_x J\le n(\ell-1)$, so
$\deg_x(v^kJ)\le L_0$ and Equation~\eqref{eq:residue-truncation}
justifies the truncation.  In a retained summand, the numerator
degree is at most $L_0+(\ell-1)L_0=\ell L_0$, and the target
degree is at most $n(\ell-1)+\ell L_0$.

The prime bound makes both degrees less than $p-1$.
The finite-field identity in Equation~\eqref{eq:field-coefficient}
therefore extracts each coefficient
by a sum over $(\mathbb F_p^*)^n$.  At a point $z$ of this grid,
its numerator is evaluated as
\[
 v(z)^kJ(z)\prod_{i=1}^n r_i(z)^{j_i}.
\]
This uses only polynomially many evaluations and ordinary modular
powering with short exponents, so lies in $\Flevel{a}$.
The grid sum and then the sum over $j$, by
Lemma~\ref{lem:counting}, place $t_k$ in $\Flevel{a+2}$.
In particular, a pair $(j,z)$ has at most
\[
 n\lceil\log_2(L_0+1)\rceil+n\lceil\log_2p\rceil
\]
bits, up to fixed padding and delimiters.

To recover characteristic coefficients from traces, recall the formal
identity for an integer matrix $A$ of dimension $D$:
\begin{equation}\label{eq:log-determinant}
 \det(I-zA)=
 \exp\left(-\sum_{k\ge1}\frac{\operatorname{tr}(A^k)}{k}z^k\right)
 \quad\text{in }\mathbb Q[[z]].
\end{equation}
Triangularizing over $\mathbb C$ reduces this to the power-series
identity $\log(1-z\alpha)=-\sum_{k\ge1}\alpha^kz^k/k$ for each
diagonal entry, and the resulting equality has rational coefficients.
This reconstruction of characteristic coefficients from residue
traces also appears in
\cite[Author version, Section~5, Equations~(5.7)--(5.9)]{CattaniDickensteinSturmfels1996}.

The matrix identity proves the required coefficient equality; the
computation uses only the indexed traces.  Using their residues,
define over $\mathbb F_p$
\begin{equation}\label{eq:truncated-exponential}
 E_v(z)=\sum_{j=0}^D\frac1{j!}
 \left(-\sum_{k=1}^D\frac{t_k}{k}z^k\right)^j.
\end{equation}
Every discarded term of the exponential identity has degree greater
than $D$.  All retained denominators are products of integers at most
$D$, and $p>D$ by Equation~\eqref{eq:char-prime-bound}.  Thus reduction
of the rational identity modulo $p$ is valid through degree $D$ and
gives
\begin{equation}\label{eq:exponential-congruence}
 E_v(z)\equiv\det(I-zM_v)\pmod{z^{D+1}}.
\end{equation}

At a given field element $z$, the inner sum in
Equation~\eqref{eq:truncated-exponential} is in $\Flevel{a+3}$.
By Lemma~\ref{lem:product}, the factorial $j!\bmod p$ is in
$\Flevel{4}$, with $0!=1$.  Modular inversion and powering take
polynomial time.  Consequently the outer sum evaluates $E_v(z)$ in
$\Flevel{\max(a+3,4)+1}$.  This entire polynomial has degree at most
$D^2<p-1$.  A final application of
Equation~\eqref{eq:field-coefficient}, in one variable, extracts a
coefficient in $\Flevel{\max(a+3,4)+2}$, which is contained in
$\Flevel{a+8}$.  The required coefficient is obtained from
\[
 [y^d]\det(yI-M_v)=[z^{D-d}]\det(I-zM_v)
                  =[z^{D-d}]E_v(z).
\]
The procedure checks primality and the numerical bounds first and
returns a fixed default outside the stated range.  Primality is
available in $\Flevel{1}$ by a divisor test, so these checks fit the
asserted level.  No step assumes squarefreeness modulo $p$.
\end{proof}

\subsection{The determinant rule for controlled families}

The preceding computation provides residues.  To obtain a controlled
family, we also need degree and coefficient-norm bounds before
specialization.  The decreasing degree in the reduction relations
supplies these bounds directly.

\begin{proposition}[Multiplication determinant rule]
\label{prop:multiplication-determinant}
Let $n\ge1$ be polynomially bounded and polynomial-time computable,
and let $\ell\ge2$ be a short polynomial-time computable integer.
Suppose that $r_1,\ldots,r_n,v$ are controlled integer polynomial
families in variables $x=(x_1,\ldots,x_n)$ and a polynomially bounded
number of parameter variables $t$, uniformly in $i$, and that
$\deg_x r_i<\ell$.  In the monomial basis of
\[
 \mathbb Z[t,x]/(x_1^\ell-r_1,\ldots,x_n^\ell-r_n),
\]
let $M_v$ denote multiplication by $v$.  Then
\begin{equation}\label{eq:controlled-characteristic}
 Q(t,y)=\det(yI-M_v)
\end{equation}
is a controlled family, monic of $y$-degree $D=\ell^n$.
If $d$ and $2^b$ are common total-degree and coefficient-norm bounds
on the $r_i,v$, and $T=d+n(\ell-1)$, valid total-degree and
log-norm bounds on $Q$ are
\begin{equation}\label{eq:determinant-size-bounds}
 d_Q=D\bigl(1+d(1+T)\bigr),\qquad
 b_Q=D^2+D\bigl(1+b(1+T)\bigr).
\end{equation}
The construction is compatible with specializing any parameters.
\end{proposition}

\begin{proof}
Lemma~\ref{lem:monomial-basis}, with $S=\mathbb Z[t]$, defines the
matrix and gives its dimension $D$ and compatibility with every
specialization.  Although $D$ need not be polynomially bounded, its
binary length is at most $1+n\lceil\log_2\ell\rceil$; hence it is
short and computable in polynomial time.

For a basis monomial $x^a$, reduce $vx^a$ to the basis by replacing
an offending factor $x_i^\ell$ by $r_i$.  In each round, perform one
such replacement on every monomial that still requires reduction,
using a fixed choice of $i$.  The initial $x$-degree is at most $T$,
and each replacement strictly decreases it, so at most $T$ rounds
suffice.  One round multiplies the coefficient-norm bound by at most
$2^b$ and increases parameter degree by at most $d$; monomials
already reduced need neither change.  Indeed, a term with coefficient
$c$ produces a polynomial of norm at most $|c|2^b$, and summing these
bounds accounts for every branch; combining like terms cannot increase
the norm.  Thus every matrix entry has
parameter degree at most $d(1+T)$ and coefficient norm at most
$2^{b(1+T)}$.

Every entry of $yI-M_v$ consequently has total degree at most
$1+d(1+T)$ and norm at most $2^{1+b(1+T)}$.
Determinant expansion and $D!\le2^{D^2}$ give precisely
Equation~\eqref{eq:determinant-size-bounds}.  The reduction argument
is used only for these bounds; the algorithm does not perform its
potentially many reductions or construct the matrix.

For evaluations, fix a prime and a point for the parameters, and
choose integer lifts of those parameter values.  These lifts are
additional short inputs; specialization gives integer polynomials in
$x$ to which Proposition~\ref{prop:char-coeff} applies.  The basis
commutes with specialization and reduction modulo $p$, so their
multiplication matrices reduce to the required matrices over the
prime field, independently of the lifts.
The controlled derivative rule in
Proposition~\ref{prop:controlled-closure} evaluates the entries of
$J=\det(\partial_jf_i)$ with a fixed increase in hierarchy level.
An ordinary $n\times n$ determinant over the field then evaluates
$J$ in polynomial bit time using these available entries.  Take
$\Omega=d$ in Proposition~\ref{prop:char-coeff} and enlarge its prime
cutoff to include those of the input families and their derivatives.
All these cutoffs are short and polynomial-time computable from the
primary word.  The proposition computes the indexed coefficients of
the specialized $Q$ by a fixed number of counting stages.  Summing
their products with $y^j$ gives the evaluation of $Q(t,y)$.

Thus $Q$ has the required available
evaluations as well as the stated short bounds, and is controlled.
\end{proof}

The rule applies equally to characteristic polynomials of coordinate
multiplication and of a polynomial objective.  In both cases it
replaces an exponentially large matrix by finite-field evaluations
with a fixed number of counting stages.  Its coefficient polynomials
are controlled by Proposition~\ref{prop:controlled-closure}, so the
rule may be used again at later, fixed stages of the construction.
